% ==============================================================================
% T0/FFGFT-Theorie: Dokument 240
% Titel: KI-Detektoren und ad-hominem-Argumentation
% Untertitel: Strukturelle Verwandtschaft alter und neuer Manipulations-Muster
% Autor: Johann Pascher
% Datum: Mai 2026
% Bezug: methodisch eigenständig, lose verwandt mit Dok 207 §10.3
% Anlass: IPI-Mailingliste, Diskussion mit Peter Austin (Mai 2026)
% ==============================================================================

\section*{Vorbemerkung}

Dieser Text geht auf eine Diskussion in der Mailingliste des
\emph{Information Physics Institute} (Mai 2026) zurück. Peter Austin
schilderte, dass seine eigenen E-Mails von KI-Detektoren regelmäßig
mit Werten über 50\,\% als KI-generiert markiert werden -- auch
dann, wenn sie in persönlicher Korrespondenz entstehen und keinerlei
KI-Unterstützung enthalten. Die Frage in der Runde war zunächst
praktisch: was tun gegen diese falschen Positiv-Treffer.

Bei näherer Betrachtung zeigt sich, dass die Frage nicht primär
technisch ist, sondern methodisch und psychologisch. KI-Detektoren
stehen nicht am Anfang eines neuen Phänomens, sondern am vorläufigen
Ende einer alten Linie: jener Argumentations-Muster, die ablehnen,
ohne sich auf die Sache einzulassen. Dieses Dokument beschreibt die
strukturelle Verwandtschaft und benennt den psychologischen
Mechanismus, der das Muster antreibt.

% ==============================================================================
\section{Der Widerspruch in der Anforderung}
% ==============================================================================

Im Zentrum der gegenwärtigen Debatte steht ein offener Widerspruch:

\begin{itemize}[leftmargin=2em,itemsep=2pt,topsep=2pt]
  \item Auf der einen Seite verlangen wir präzise, fehlerfreie,
        klar strukturierte Texte. Das ist seit jeher der Standard
        wissenschaftlicher und fachlicher Kommunikation.
  \item Auf der anderen Seite gilt nun \emph{die Abwesenheit von
        Fehlern} selbst als verdächtig. Ein Text mit niedriger
        Perplexity, gleichmäßiger Satzstruktur und sauberen
        Übergängen wird von Detektoren als KI markiert.
\end{itemize}

Wörtlich genommen ist das eine \textbf{Aufforderung, schlechter zu
schreiben}: genug Roughness produzieren, um nicht aufzufallen, genug
Variation, um nicht als Maschine zu gelten. Das kann kein
methodischer Standard sein.

Die Belege für diese Schieflage sind klar dokumentiert. Die
Stanford-Studie von Weixin Liang, Mert Yuksekgonul, Yining Mao,
Eric Wu und James Zou \emph{GPT detectors are biased against
non-native English writers} (Patterns, Vol.~4, Nr.~7, Artikel
100779, 14. Juli 2023; arXiv:2304.02819) zeigte: Sieben weit
verbreitete GPT-Detektoren klassifizierten \emph{über die Hälfte}
der TOEFL-Aufsätze nicht-muttersprachlicher Englisch-Lerner als
\enquote{KI-generiert}, während dieselben Detektoren bei
US-amerikanischen Achtklässler-Aufsätzen nahezu perfekte Genauigkeit
erreichten. Der Grund: nicht-muttersprachliche Texte haben
typischerweise eine niedrigere Perplexity, weil das Vokabular
eingeschränkter und die Satzstrukturen einfacher gewählt werden --
genau dieses Merkmal teilen sie mit maschinell generiertem Text.

OpenAI selbst -- der Hersteller von ChatGPT -- veröffentlichte im
Januar~2023 einen \emph{AI Classifier}, der genau dieses Problem
lösen sollte. Sechs Monate später, am 20.~Juli~2023, wurde das
Werkzeug eingestellt. Der offizielle Grund laut OpenAI: nur
\emph{26\,\%} der tatsächlich KI-generierten Texte wurden korrekt
erkannt, und in \emph{9\,\%} der Fälle wurden menschliche Texte
fälschlich als KI markiert. Bei einer Trefferquote unter 50\,\%
ist ein Detektor schlechter als ein Münzwurf.

Wer als Techniker, Wissenschaftler oder Nicht-Muttersprachler viel
überarbeitet, schreibt notwendigerweise klarer und strukturierter
als Gelegenheitsschreiber. Genau dieses Schreiben wird systematisch
geflaggt.

% ==============================================================================
\section{Substanz statt Stil}
% ==============================================================================

Die eigentliche Unterscheidung zwischen menschlich und maschinell
generiertem Inhalt liegt nicht im Stil, sondern in der Substanz.

\textbf{KI-Texte sind charakterisiert nicht durch ihre Form, sondern
durch eine spezifische inhaltliche Schwäche:} sie sind oft
schlüssig klingende Argumentationen, die auf
\emph{halluzinierten Fakten}, \emph{erfundenen Zitaten} und
\emph{selbstsicher vorgetragenen, aber inhaltlich leeren
Herleitungen} aufgebaut sind. Die sprachliche Politur ist real;
die Substanz ist erfunden.

Ein kritischer Leser, der den Gegenstand versteht, erkennt diese
Schwäche -- nicht an Oberflächenmerkmalen, sondern daran, dass die
Behauptungen bei genauerer Prüfung nicht halten. Eine statistische
Maschine kann diesen Unterschied prinzipiell nicht sehen, denn er
ist kein Oberflächenmerkmal.

\textbf{Folge:} Wer Inhalt prüfen will, muss Inhalt verstehen
können. Detektoren können das nicht. Sie können nur Symptome
zählen.

% ==============================================================================
\section{Der Checker-Nutzer als eigentlicher Fehlanwender}
% ==============================================================================

Hier zeigt sich eine Inversion, die in der Debatte selten ausgesprochen
wird. Wer einen KI-Detektor benutzt, statt selbst zu lesen,
\textbf{delegiert seine eigene Urteilskraft an einen Algorithmus}.

Damit ist der Checker-Nutzer tiefer in der KI-Abhängigkeit als der
Autor, der KI bewusst als Werkzeug nutzt:

\begin{center}
\renewcommand{\arraystretch}{1.4}
{\small
\begin{tabular}{>{\raggedright\arraybackslash}p{3.4cm}>{\raggedright\arraybackslash}p{3.6cm}>{\raggedright\arraybackslash}p{3.4cm}}
\toprule
\textbf{Rolle} & \textbf{Was tut die Person} & \textbf{Wer hat das Urteil} \\
\midrule
Autor mit KI-Werkzeug & Substanz selbst, Sprache mit KI-Hilfe & Mensch \\[3pt]
Reviewer mit Detektor & Score ablesen, Urteil von Maschine übernehmen & Maschine \\
\bottomrule
\end{tabular}
}
\renewcommand{\arraystretch}{1.0}
\end{center}

Der Vorwurf \enquote{Sie haben KI benutzt} kommt also häufig von
genau der Person, die ihre eigene kritische Funktion bereits an KI
abgegeben hat. Sie wirft dem Autor vor, was sie selbst tut -- ohne
Bewusstsein, dass sie es tut.

% ==============================================================================
\section{KI als Filter, nicht als Orakel}
% ==============================================================================

Eine wichtige Differenzierung ist nötig. Auch sorgfältige
Wissensarbeiter nutzen KI -- für Zusammenfassungen, für
Sprachformalisierung, für Vorsortierung großer Textmengen. Das ist
nicht das Problem. Das Problem ist \emph{die Art} der Nutzung.

\textbf{KI als Filter (Triage):}
\begin{itemize}[leftmargin=2em,itemsep=2pt,topsep=2pt]
  \item KI liest große Textmengen vor und fasst zusammen.
  \item Der Mensch liest die Zusammenfassung und entscheidet:
        passt das, oder muss ich selbst genauer hinsehen?
  \item Bei Unklarheit oder Auffälligkeit: der Mensch geht selbst
        ins Detail.
  \item Das Endurteil liegt immer beim Menschen.
\end{itemize}

\textbf{KI als Orakel (Delegation):}
\begin{itemize}[leftmargin=2em,itemsep=2pt,topsep=2pt]
  \item KI gibt einen Score, eine Klassifikation, eine Wahrheit.
  \item Der Score wird übernommen.
  \item Keine eigene Prüfung der Substanz.
  \item Keine Bereitschaft, bei Auffälligkeiten selbst zu lesen.
\end{itemize}

Die KI-Nutzung als \textbf{Filter} ist methodisch unproblematisch
-- sie entspricht der seit jeher üblichen Praxis, größere
Textmengen durch Abstracts, Bibliographien oder Assistenten
vorsortieren zu lassen. Die KI-Nutzung als \textbf{Orakel} ist
methodisch fatal: sie ersetzt das Urteil, statt es zu unterstützen.

Detektoren sind das Paradebeispiel der Orakel-Nutzung. Sie liefern
eine Zahl, die akzeptiert oder ignoriert werden kann -- aber keine
Substanz, die geprüft werden könnte.

% ==============================================================================
\section{Strukturelle Unmöglichkeit der Inhaltsbeurteilung}
% ==============================================================================

Hier liegt die technische Pointe, die der ganzen Detektor-Konstruktion
den Boden entzieht.

Ein KI-Detektor behauptet implizit: \enquote{Ich kann erkennen, ob
ein Text von einer KI stammt.} Dafür müsste er etwas können, was
die KI selbst nicht kann:

\begin{itemize}[leftmargin=2em,itemsep=2pt,topsep=2pt]
  \item KI weiß nicht, ob eine Behauptung wahr ist.
  \item KI weiß nicht, ob eine Argumentation hält.
  \item KI erkennt nicht ihre eigenen Halluzinationen.
\end{itemize}

\textbf{Wenn die KI selbst inhaltlich überfordert ist, kann ein
Detektor -- der ebenfalls eine KI ist -- erst recht nicht
inhaltlich urteilen.} Er kann nur Oberflächenmerkmale messen:
Wortwahrscheinlichkeiten (Perplexity), Satzlängen-Variation
(Burstiness), Häufigkeit bestimmter Übergangswörter.

Die Konstruktion löst sich logisch selbst auf:

\begin{quote}
Wer Inhalt prüfen will, muss Inhalt verstehen können. Detektoren
sind KI. KI kann keinen Inhalt verstehen. Also können Detektoren
nicht prüfen, was sie zu prüfen vorgeben.
\end{quote}

Ein Detektor, der wirklich Inhalt prüfen könnte, wäre keine KI
mehr -- er wäre ein \emph{kritisch denkender Mensch mit Sachverstand}.
Genau das soll der Detektor aber ersetzen. Der Anspruch ist
selbstwidersprüchlich.

\subsection{Eine Analogie}

Niemand würde jemandem glauben, der sagt: \enquote{Ich kann
erkennen, ob ein Mathebuch von einem inkompetenten Autor stammt --
ohne dass ich selbst Mathematik kann. Ich messe nur die Häufigkeit
bestimmter Wörter.} Aber genau das tun KI-Detektoren: Sie
behaupten, \textbf{Substanz} durch \textbf{Form-Statistik} zu
beurteilen.

Das ist derselbe Kategorienfehler wie bei Phrenologie (Schädelform
gleich Charakter) oder Lügendetektor (Schweißproduktion gleich
Wahrheit): \emph{Messung von Oberflächen-Korrelaten, dargestellt
als Messung der eigentlichen Größe.}

% ==============================================================================
\section{Verwandtschaft mit ad-hominem-Argumentation}
% ==============================================================================

Bis hierhin war die Argumentation technisch. Nun kommt die
strukturelle Pointe: KI-Detektoren sind kein neues Phänomen.

\textbf{Sie machen mit Texten dasselbe, was ad-hominem-Argumentation
mit Personen macht.}

\begin{center}
\renewcommand{\arraystretch}{1.3}
{\small
\begin{tabular}{>{\raggedright\arraybackslash}p{5.0cm}>{\raggedright\arraybackslash}p{5.6cm}}
\toprule
\textbf{Ad-hominem-Argumentation} & \textbf{KI-Detektor} \\
\midrule
greift den Sprecher an, nicht das Argument & bewertet die Form, nicht den Inhalt \\[2pt]
wer hat das gesagt? & wer (oder was) hat das geschrieben? \\[2pt]
umgeht die inhaltliche Auseinandersetzung & umgeht die inhaltliche Prüfung \\[2pt]
ist bequem -- kein Verstehen nötig & ist bequem -- kein Lesen nötig \\[2pt]
ersetzt Denken durch Verdacht & ersetzt Lesen durch Score \\
\bottomrule
\end{tabular}
}
\renewcommand{\arraystretch}{1.0}
\end{center}

Beides folgt demselben Muster: \textbf{Verlagerung der Bewertung
von der Sache auf eine Eigenschaft der Quelle.} Statt zu prüfen,
ob etwas stimmt, wird gefragt, woher es kommt -- und aus der
Herkunft auf die Gültigkeit geschlossen.

In der Logik heißt das seit Jahrhunderten \emph{genetischer
Fehlschluss}. Aristoteles hat ad-hominem-Argumente bereits in den
\emph{Sophistischen Widerlegungen} als Fehlschluss beschrieben.
Arthur Schopenhauer behandelt das verwandte \emph{argumentum ad
personam} in seiner \emph{Eristischen Dialektik} als Scheinargument,
das den Inhalt umgeht. Der Mechanismus ist 2400 Jahre alt; die
technische Vollendung in Software-Form ist neu.

% ==============================================================================
\section{Kognitive Dissonanz als Antrieb des Musters}
% ==============================================================================

Warum funktioniert dieses Muster so zuverlässig? Die Antwort liefert
Leon Festingers \emph{Theory of Cognitive Dissonance} (Stanford
University Press, 1957), eine der einflussreichsten Arbeiten der
Sozialpsychologie des 20.~Jahrhunderts.

\subsection{Festingers Grundbeobachtung}

Festinger entwickelte die Theorie aus einer ungewöhnlichen
empirischen Beobachtung. In \emph{When Prophecy Fails} (gemeinsam
mit Henry W.~Riecken und Stanley Schachter, University of Minnesota
Press, 1956) dokumentiert er das Verhalten einer Weltuntergangs-
Sekte, die einen konkreten Termin für das Ende der Welt vorhergesagt
hatte. Als der Termin verstrich, ohne dass etwas geschah,
\emph{verstärkte} sich der Glaube der Mitglieder, statt zu
zerfallen. Sie deuteten das Ausbleiben des Untergangs als Beweis,
dass ihre Glaubensgemeinschaft die Welt gerettet hatte.

Festingers Schlussfolgerung: Wenn zwei Überzeugungen unvereinbar
sind, entsteht psychischer Druck (\emph{Dissonanz}). Dieser Druck
führt nicht zwingend zur Korrektur der falschen Überzeugung -- oft
führt er zur \emph{Umdeutung der Realität}, so dass die ursprüngliche
Position bestätigt bleibt.

\subsection{Anwendung auf den Detektor-Kontext}

In der KI-Detektor-Situation entsteht typischerweise diese
Dissonanz beim Reviewer:

\begin{center}
\renewcommand{\arraystretch}{1.3}
{\small
\begin{tabular}{>{\raggedright\arraybackslash}p{5.0cm}>{\raggedright\arraybackslash}p{5.6cm}}
\toprule
\textbf{Überzeugung A} & \textbf{Überzeugung B (auslösend)} \\
\midrule
\enquote{Ich bin ein kompetenter Reviewer} & \enquote{Vor mir liegt ein Text klarer als das, was ich schreiben würde} \\[2pt]
\enquote{Ich gehöre zu denen, die kritisch denken} & \enquote{Hier denkt offenbar jemand klarer als ich} \\[2pt]
\enquote{KI macht Menschen dümmer} & \enquote{Aber dieser Text wirkt nicht dumm, sondern klüger als üblich} \\
\bottomrule
\end{tabular}
}
\renewcommand{\arraystretch}{1.0}
\end{center}

Die Dissonanz ist unangenehm. Es gibt zwei Wege, sie aufzulösen:

\textbf{Weg~1: Inhaltliche Auseinandersetzung} (anstrengend).\\
Die eigene Position überdenken; anerkennen, dass der andere etwas
Gutes geleistet hat; eigene Standards nachjustieren; bei Bedarf die
eigene Position revidieren.

\textbf{Weg~2: Quellenabwertung} (bequem).\\
Den Text als nicht-authentisch markieren -- \enquote{das ist ja gar
nicht von ihm, das ist KI}. Die Selbstwahrnehmung bleibt intakt.
Keine eigene Revision nötig.

Der KI-Detektor liefert die \emph{technische Rechtfertigung} für
Weg~2. Er erlaubt es, den unbequemen Text zu disqualifizieren, ohne
sich mit ihm auseinandersetzen zu müssen. Der Score wirkt
wissenschaftlich, objektiv, neutral -- und besorgt genau das, was
der Reviewer braucht, um seine Dissonanz aufzulösen, ohne sich
selbst zu hinterfragen.

\subsection{Drei Funktionen in einem}

Die Dissonanz-Reduktion durch Quellenabwertung erfüllt
gleichzeitig drei Funktionen:

\begin{itemize}[leftmargin=2em,itemsep=2pt,topsep=2pt]
  \item \textbf{Selbstwertschutz:} \enquote{Ich bin nicht weniger
        kompetent -- der hat nur gemogelt.}
  \item \textbf{Anstrengungsvermeidung:} \enquote{Ich muss mich
        nicht damit auseinandersetzen.}
  \item \textbf{Soziale Positionierung:} \enquote{Ich gehöre zu
        den Kritischen, die KI-Mogelei erkennen.}
\end{itemize}

Alle drei Funktionen sind \emph{billig} zu haben -- ein Klick auf
einen Detektor. Die ehrliche Alternative -- inhaltliche Prüfung --
ist \emph{teuer}: Zeit, Konzentration, das Risiko der eigenen
Revision.

\subsection{Die Verschärfung durch Pseudo-Objektivität}

Klassische ad-hominem-Argumente können vom Adressaten direkt
zurückgewiesen werden -- man kann antworten, kontextualisieren,
sich wehren. KI-Detektoren haben drei zusätzliche Eigenschaften,
die die Dissonanz-Reduktion noch effektiver machen:

\begin{enumerate}[leftmargin=2em,itemsep=2pt,topsep=2pt]
  \item \textbf{Pseudo-Objektivität:} Die Zahl wirkt wissenschaftlich,
        nicht persönlich. Der Reviewer kann sagen: \enquote{Ich
        greife Sie ja nicht an -- der Algorithmus sagt das.}
  \item \textbf{Verantwortungs-Auslagerung:} Niemand fühlt sich
        persönlich verantwortlich für das Urteil. \enquote{Der
        Computer hat entschieden.}
  \item \textbf{Unwiderlegbarkeit:} Der Adressat kann \emph{nicht
        beweisen}, dass er die KI nicht benutzt hat. Negation ist
        nicht beweisbar.
\end{enumerate}

Diese drei Eigenschaften machen den KI-Detektor zu einem
\emph{idealen Werkzeug zur Dissonanz-Reduktion}: maximale Wirkung
bei minimaler eigener Investition und minimaler Verantwortung.

\subsection{Eine bittere Pointe}

Daraus folgt eine paradoxe Konsequenz: Je \textbf{klarer und
kompetenter} ein Text ist, desto \textbf{größer die Dissonanz} für
einen weniger souveränen Leser -- und desto \textbf{stärker das
Motiv}, den Text als KI abzustempeln. Die Detektoren bestrafen also
genau die Kompetenz, die sie zu schützen vorgeben. Nicht zufällig,
sondern aus psychologischer Funktion.

% ==============================================================================
\section{Die ganze Familie rhetorischer Manipulations-Muster}
% ==============================================================================

Die Verwandtschaft ist nicht auf den ad-hominem-Fehlschluss
beschränkt. KI-Detektoren reproduzieren tatsächlich \emph{eine ganze
Familie} klassischer Manipulations-Techniken -- jeweils in
technischer Form.

\subsection{Tone Policing}
\textbf{Mensch:} \enquote{Sie sind aber emotional -- das kann man
so nicht ernst nehmen.}\\
\textbf{Detektor:} \enquote{Sie schreiben aber sehr glatt -- das
kann man so nicht ernst nehmen.}

Der Stil wird zum Disqualifikator, unabhängig vom Inhalt.

\subsection{Strohmann-Argument}
Der \emph{Mensch} verzerrt die Position des Gegners zu einer
einfacheren, leichter zu widerlegenden Karikatur. Der \emph{Detektor}
reduziert den Text auf statistische Oberflächenmerkmale, die mit
dem eigentlichen Inhalt nichts zu tun haben. In beiden Fällen wird
eine vereinfachte Karikatur bewertet, nicht das Original.

\subsection{Beweislastumkehr}
\textbf{Mensch:} \enquote{Beweis mir, dass Du das selber gedacht
hast.}\\
\textbf{Detektor:} \enquote{Beweise mir, dass Du keine KI benutzt
hast.}

In beiden Fällen wird das Gegenteil als Default angenommen. Der
Adressat muss seine Unschuld beweisen -- was strukturell unmöglich
ist, weil \emph{Nicht-Existenz nicht beweisbar} ist.

\subsection{Argumentum ad verecundiam}
\textbf{Mensch:} \enquote{Das hat Professor~X gesagt, also muss es
stimmen.}\\
\textbf{Detektor:} \enquote{Der Algorithmus sagt 87\,\%, also muss
es stimmen.}

Eine Autorität wird über den Inhalt gestellt. Bei KI-Detektoren ist
die Autorität eine Maschine -- was die Inversion verschärft.

\subsection{Argumentum ad populum}
\textbf{Mensch:} \enquote{Alle wissen doch, dass\ldots}\\
\textbf{Detektor:} \enquote{Das sieht aus wie typischer KI-Stil,
alle Modelle erkennen das Muster.}

Konsens (oder gelerntes Muster) ersetzt Argument.

\subsection{Gaslighting}
\textbf{Mensch:} \enquote{Sie überschätzen sich -- das kann nicht
von Ihnen kommen.}\\
\textbf{Detektor:} \enquote{Der Score zeigt: das kann nicht von
Ihnen kommen.}

Dem Adressaten wird die eigene Autorenschaft abgesprochen. Er soll
an sich selbst zweifeln. Das ist die destruktivste Form -- sie
zielt nicht auf die Aussage, sondern auf die Person selbst.

% ==============================================================================
\section{Die historische Linie}
% ==============================================================================

Die hier beschriebenen Muster sind nicht neu. Sie sind in jeder
Epoche dokumentiert:

\begin{itemize}[leftmargin=2em,itemsep=2pt,topsep=2pt]
  \item \textbf{Antike Rhetorik:} Aristoteles hat ad-hominem-
        Argumente in den \emph{Sophistischen Widerlegungen} als
        Fehlschlüsse beschrieben.
  \item \textbf{Mittelalter:} \emph{Argumentum ad auctoritatem}
        (Berufung auf Autoritäten statt eigene Prüfung) als
        scholastische Argumentationsform.
  \item \textbf{Inquisition:} Verdachts-Logik -- wer so denkt, ist
        verdächtig.
  \item \textbf{Totalitarismus des 20.~Jahrhunderts:} \enquote{Wer
        so spricht, ist Feind.}
  \item \textbf{Gegenwart:} KI-Detektoren -- \enquote{Wer so
        schreibt, ist nicht echt.}
\end{itemize}

Es ist immer dasselbe Muster: Wenn die inhaltliche Auseinandersetzung
schwierig wird, weicht man auf die Quellenkritik aus. Die
\textbf{historische Konstanz} ist die wichtigste Beobachtung dieses
Dokuments -- sie zeigt, dass das eigentliche Problem nicht
technologisch, sondern \emph{anthropologisch} ist.

% ==============================================================================
\section{Was tatsächlich helfen würde}
% ==============================================================================

Wenn die Detektoren nicht helfen, sondern das Problem verschärfen,
was hilft dann?

\textbf{Substantielles Lesen.} Ein kompetenter Leser, der den
Gegenstand versteht, erkennt Halluzinationen, weil er die Substanz
versteht. Er braucht keinen Score. Er braucht Bereitschaft, sich
auf den Inhalt einzulassen.

\textbf{Reproduzierbare Verifikation.} Wo der Inhalt prüfbar ist,
kann er geprüft werden -- durch Code, Daten, Skripte, die man
selbst ausführen kann. Das ist keine Detektor-Funktion; das ist
\emph{Inhaltsfunktion}.

\textbf{Provenance-Dokumentation.} Drafts, Aufzeichnungen,
Revisionsverlauf. Wer transparent macht, wie er gearbeitet hat,
verschiebt die Frage von \enquote{Beweis, dass Du keine KI benutzt
hast} zu \enquote{hier ist der Arbeitsverlauf, schau selbst}. Die
zweite Frage ist beantwortbar.

\textbf{Offene Disziplin.} KI bewusst als Werkzeug nutzen, nicht
als Orakel. Eigene Position als die eigene markieren. Die KI-Hilfe
offen ausweisen, statt sie zu verstecken.

Keiner dieser Punkte braucht einen Detektor. Alle vier brauchen
\emph{kompetente, ehrliche Menschen, die bereit sind, die Arbeit
des Verstehens zu tragen}.

% ==============================================================================
\section{Diagnose}
% ==============================================================================

Die KI-Detektor-Debatte ist nicht primär technisch. Sie ist die
sichtbare Spitze einer kulturellen Tendenz, die das anstrengende
Geschäft des Verstehens durch die bequeme Praxis der Quellenkritik
ersetzt.

In einem Satz:

\begin{quote}
\textbf{Ad-hominem-Argumentation ist KI-Detektor für Menschen. KI-
Detektor ist ad-hominem-Argumentation für Texte. Beides ist die
Verweigerung der Sache, getarnt als Sorge um die Quelle.}
\end{quote}

Der psychologische Antrieb -- kognitive Dissonanz und ihre Auflösung
durch Quellenabwertung -- ist seit Festinger 1956/1957 wohl
dokumentiert. Die historische Konstanz des Musters reicht bis
Aristoteles zurück. Die technische Vollendung in Software-Form ist
neu, das zugrundeliegende Problem ist alt.

Wer einem KI-Checker mehr vertraut als seinem eigenen kritischen
Urteil, hat aufgehört, ein Reviewer zu sein -- und ist selbst zum
Anhängsel einer Maschine geworden. Das ist der eigentliche
Fehlanwender von KI.

% ==============================================================================
\section{Quellen}
% ==============================================================================

\begin{itemize}[leftmargin=2em,itemsep=2pt,topsep=2pt]
  \item Liang, W., Yuksekgonul, M., Mao, Y., Wu, E., Zou, J.
        (2023). \emph{GPT detectors are biased against non-native
        English writers}. Patterns Vol.~4, Nr.~7, Artikel 100779.
        DOI:~10.1016/j.patter.2023.100779; arXiv:2304.02819.
  \item OpenAI (2023). \emph{New AI classifier for indicating
        AI-written text}. Blog-Beitrag vom 31.~Januar 2023, ergänzt
        am 20.~Juli 2023 um den Hinweis auf die Einstellung des
        Werkzeugs wegen niedriger Trefferquote.
  \item Festinger, L. (1957). \emph{A Theory of Cognitive
        Dissonance}. Stanford University Press.
  \item Festinger, L., Riecken, H.~W., Schachter, S. (1956).
        \emph{When Prophecy Fails: A Social and Psychological
        Study of a Modern Group that Predicted the Destruction of
        the World}. University of Minnesota Press.
  \item Aristoteles. \emph{Sophistische Widerlegungen} (Sophistikoi
        Elenchoi). Klassische Darstellung der rhetorischen Fehlschlüsse.
  \item Schopenhauer, A. \emph{Eristische Dialektik} (postum
        veröffentlicht). Behandlung des \emph{argumentum ad personam}.
\end{itemize}

\bigskip

\textbf{Anmerkung zur Entstehung dieses Dokuments.} Der Text wurde
im Mai 2026 verfasst, ausgelöst durch eine Diskussion auf der
Mailingliste des \emph{Information Physics Institute} mit Peter
Austin. Die Argumentation ist mein eigenes Denken. Die sprachliche
Formalisierung erfolgte mit KI-Unterstützung -- in genau der Weise,
die in §5 als \emph{KI als Filter, nicht als Orakel} beschrieben
ist. Eine offene Disziplin: Substanz vom Autor, Sprache mit
Werkzeug-Unterstützung, Urteil beim Menschen.
